\documentclass[12pt]{article}

% --------------------
% PAQUETES BÁSICOS
% --------------------
\usepackage[spanish,es-noshorthands]{babel}
\usepackage[utf8]{inputenc}      % Si usas pdflatex
\usepackage[T1]{fontenc}
\usepackage{lmodern}

\usepackage{amsmath, amssymb, amsthm}
\usepackage{geometry}
\usepackage{booktabs}   
\usepackage{setspace}
\usepackage{graphicx}
\usepackage{tikz}
\usepackage{tikz-cd}
\usepackage[
  colorlinks=true,
  linkcolor=blue,
  citecolor=blue,
  urlcolor=red
]{hyperref}

\newtheorem{thm}{Teorema}[section]
\newtheorem{dfn}[thm]{Definición}
\newtheorem{lem}[thm]{Lema}
\newtheorem{ej}[thm]{Ejemplo}
\newtheorem{cor}[thm]{Corolario}
\newtheorem{obs}[thm]{Observación}
\newtheorem{con}[thm]{Conjetura}
\newtheorem{prop}[thm]{Proposición}

\DeclareMathOperator{\Frm}{Frm}
\DeclareMathOperator{\Ord}{Ord}
\DeclareMathOperator{\Top}{Top}
\DeclareMathOperator{\pt}{pt}
\DeclareMathOperator{\id}{id}
\DeclareMathOperator{\tp}{tp}
\DeclareMathOperator{\fd}{fd}
% --------------------
% CONFIGURACIÓN
% --------------------
\geometry{margin=2.5cm}
\onehalfspacing

% --------------------
% DATOS DEL DOCUMENTO
% --------------------
\title{Reporte mensual de actividades}
\date{1-marzo-2026 a 31-marzo-2026}
\author{Juan Carlos Monter Cortés}

% --------------------
% INICIO DEL DOCUMENTO
% --------------------
\begin{document}

\maketitle
% \tableofcontents

\section{Actividades realizadas}

En \cite{sexton2006ordinal, simmons2004vietoris} prueban algunas propiedades que satisface un espacio topológico cuando este es apilado o fuertemente 
apilado. Recordemos que en nuestro trabajo hemos definido los análogos de apilado y fuertemente apilado de la siguiente manera.

\begin{dfn}\label{def:apilado-marcos}
Sea $A\in \Frm$.
\begin{enumerate}
  \item Diremos que $A$ es fuertemente apilado si $\mathcal{U}_*[Q']\mathcal{U}=v_F$.
  \item Diremos que $A$ es apilado $(\mathcal{U}_*([Q'])\mathcal{U})(0)=v_F(0)$.
\end{enumerate}
\end{dfn}

Una pregunta natural que nos podríamos realizar con respecto a las definiciones anteriores es si los marcos apilados y fuertemente 
apilados se comportan de manera similar a las definiciones originales en espacio.\\

Por ejemplo, si un espacio $S$ es fuertemente apilado y $T_1$, para la par ($\nabla\in \mathcal{O}S^\wedge, Q\in \mathcal{Q}S$) se cumple que
\begin{equation}\label{v=w}
v_\nabla=w_\nabla,
\end{equation}
donde $v_\nabla$ y $w_\nabla$ son el menor y el mayor elemento del intervalo de admisibilidad correspondiente al filtro $\nabla$, respectivamente.\\

De manera general, si consideramos $A\in \Frm$, $F\in A^\wedge$ y $S=\pt A$, el Teorema de Hofmann-Mislove proporciona los conjuntos 
\[
M=\{m\in A\setminus F\mid a\leq m\}\quad \text{y}\quad Q=S\setminus F
\]
y el intervalo de admisibilidad asociado a $F$ tiene la forma
\[
v_F\leq [Q']\leq w_F=[M'].
\]

Si $S$ es un espacio $T_1$, se cumple que $M=Q$ y el intervalo de admisibilidad es
\[
v_F\leq [Q']=[M'].
\]

Si $S$ es fuertemente apilado, se satiface que $v_F=[Q']$ y el intervalo de admisibilidad es
\[
v_F=[Q']\leq [M'].
\]

Por lo tanto, si $S$ es un espacio $T_1$ y fuertemente apilado, se cumple que $v_F=[Q']=[M']$, es decir, $v_F=w_F$.\\

El análogo de la afirmación anterior para marcos es el siguiente. 

\begin{prop}\label{fapiladoT1}
Sean $A\in \Frm$ y $(F\in A^\wedge,Q \in \mathcal{Q}S, \nabla\in \mathcal{O}S^\wedge)$.
Si $A$ es $T_1$ y fuertemente apilado, entonces
\[
v_F=w_F.
\]
\end{prop}

\begin{proof}

\end{proof}

El Teorema 7.11 de \cite{simmons2004vietoris} sugiere probar el siguiente resultado.
\begin{prop}\label{focales}
Un marco $A$ es $T_1$ y apilado si y solo si $v_F(0)=w_F(0)$.
\end{prop}

\begin{proof}

\end{proof}

El Teorema 7.11 antes mencionado tiene relación con la construcción de la $V$-modificación de un espacio. Por tal motivo,
resulta de nuestro interés probar ahora que ocurre en el contexto de la $V$-modificación de un marco. Por ejemplo, para
\cite[Teorema 7.8]{sexton2006ordinal}, ¿podemos tener su análogo usando las respectivas construcciones en $\Frm$?\\

Adicionalmente, el Teorema 6.2 a) de \cite{sexton2006ordinal} dice que si un espacio 
si un espacio $S$ es $T_1$ y apilado, entonces $S$ es sobrio. Trasladando esta afirmación al contexto de marcos tenemos lo siguiente:\\

\textbf{Afirmación:} Si un marco $A$ es $T_1$ y apilado, entonces $A$ es espacial.\\

Dentro de las actividades realizadas durante este mes se encuentra el querer demostrar esta afirmación. 
Para esto, hemos partido del hecho de que un marco es espacial si y solo si el morfismo reflexión espacial 
$\mathcal{U}$ es inyectivo.\\

Abordando lo anterior a manera de contradicción, consideramos $A\in\Frm$ un marco $T_1$ y apilado, y suponemos que $\mathcal{U}$ no es inyectivo. 
Así, existen $a,b\in A$ con $a\nleq b$ tal que $\mathcal{U}(a)=\mathcal{U}(b)$, es decir,
\[
[\forall p\in \pt A](p\in \mathcal{U}(a) \iff p\in \mathcal{U}(b)).
\]

Alternativamente, si $M_a=\mathcal{U}(a)'$ y $M_b=\mathcal{U}(b)'$, tenemos que
\[
[\forall p\in \pt A](p\in M_a \iff p\in M_b).
\]

con esto, la prueba se reduce a exhibir un punto $p\in \pt A$ tal que
\[
a\nleq p \text{ y } b\leq p\quad o \quad a\leq p \text{ y } b\nleq p.
\]

Por la correspondencia entre puntos y filtros completamente primos, es suficiente mostrar que existe un filtro completamente primo $F$ tal que
\[
a\in F \text{ y } b\notin F \quad o \quad a\notin F \text{ y } b\in F.
\]

El enfoque anterior también se puede abordar a través del Teorema de Hofmann-Mislove. Si consideramos un par HM ($F\in A^\wedge, Q\in \mathcal{Q}S$), tenemos que 
\[
\begin{split}
v_F(0)=(\mathcal{U}\circ [Q']\circ \mathcal{U})(0)&=\bigvee \{x\in A\mid [\forall q \in \pt A](q\in Q\Rightarrow x\leq q)\}.
\end{split}
\]

De manera particular, si consideramos $p\in \pt A$ tenemos el par HM ($F_p, \uparrow p$), donde 
\[
F_p=\{x\in A\mid [q\in \pt a](q\leq p\Rightarrow x\nleq q)\}. 
\]

Como $A$ es $T_1$, se cumple que todo punto es máximo. Así,
\[
F_p=\{x\in A\mid x\leq p\}.
\]

\bibliographystyle{amsalpha}
\bibliography{research2}


\end{document}